
\section{Introduction}

This paper takes a further step towards improved computational schemes to high-dimensional problems described by time-dependent Hamilton-Jacobi (HJ) partial differential equations (PDEs)~\citet{EvansPDEBook}. These PDEs are increasingly employed in the safety verification~\citet{KerianneHobbs} of cyberphysical systems. We work within a viscous HJ reachability ramification~\citet{LygerosReachability, Mitchell2005, MoluxLevPyCDC, LevelSetTOMS}  and establish our results with globally optimal solutions to reachability problems in nonlinear settings. In emphasis, we consider  first-order nonlinear scalar HJ PDEs of the form % initial-value problem (IVP)  
%
%\begin{subequations}
	\begin{align}
		\partial_t\valuefunc + \hamfunc(t; \state, \partial_x \valuefunc) = 0 \text{   in   } \ren \times (0, T], \quad 
		\valuefunc(0; \state) = \bm{g}(\state) \text{   on   } \ren \times \{t=0\}
		\tag{HJ-IVP}
		\label{eq:ivp}
	\end{align}
%\end{subequations}
%
where $\partial_t\valuefunc$ denotes the partial derivative of the vector field $\valuefunc(\state,t)$ with respect to time, $t$; the Hamiltonian %$H \in C(\ren)$ \ie 
$\hamfunc(\state,t)$ is continuously defined on $\ren$; $\bm{g}(\state)$ is bounded and uniformly continuous (BUC) in $\ren$; and $\partial_x \valuefunc$ is the spatial gradient of $\valuefunc(\state,t)$. 

While a unique $C^2$ solution $\valuefunc(\state,t)$ exists uniformly on $0 \le t < T$ to \eqref{eq:ivp} for which $\lim_{t \rightarrow T} \valuefunc(\state,t)$  (when $\hamfunc$ and $\bm{g}$ are smooth and $\valuefunc_0(\state,t)$ is compactly supported), $\hamfunc(\state,t)$ lacks continuous derivatives in this limit so that $\partial_x \valuefunc$ \textit{vanishes} at $t=T$. %In this article, we only care for \textit{approximate solutions}. 
However, one may find Lipschitz continuous functions $\valuefunc(\state,t)$ on $\ren \times (0, T]$ which are smooth \textit{almost everywhere} (a.e.) and satisfy \eqref{eq:ivp} in a ``generalized sense". Suppose that we fix a viscosity term $\delta>0$, Crandall and Lions~\citet{CrandallTwoApprox} showed that $\valuefunc^\delta$ is the \textit{vanishing viscosity} solution of 
%
%\begin{subequations}
	\begin{align}
		\partial_t\valuefunc^\delta + \hamfunc(t; \state, \partial_x \valuefunc^\delta) = \dfrac{\delta}{2} \Delta \valuefunc^\delta \text{ in } \ren \times (0, T]; \,
		\valuefunc^\delta(0; \state) = \bm{g}(\state) \text{ on } \ren \times \{t=0\}
		\label{eq:ivp-viscous}
		\tag{HJ-Visc}
	\end{align}
%
with Laplacian $\Delta \valuefunc^\delta$. For a constant $k>0$, uniform convergence of $\valuefunc^\delta \rightarrow \valuefunc$ exists in the limit $\delta \rightarrow 0^+$, \ie
%
\begin{align}
	\sup_{t \in (0, T)} \sup_{x\in \ren} \mid \valuefunc(t; \state) - \valuefunc^\delta(t; \state)\mid \le k\sqrt{\delta}.
\end{align}
%
% Equation \eqref{eq:ivp-viscous} establishes .

%Equation \eqref{eq:ivp-viscous} models conservation laws in a single spatial dimension~\citet{OsherShuENO}. 
In $\inf$-$\sup$ or $\sup$-$\inf$ optimal control problems~\citet{LygerosReachability}, the Hamiltonian of \eqref{eq:ivp-viscous} is related to the \textit{backward} reachable set~\citet{Mitchell2005} of a dynamical system. Reachable systems are those where one can compute all states that can be reached from an initial condition in \textit{a finite number of steps}. %Mitchell~\citet{Mitchell2005} related levelset schemes to reachability analysis in optimal control, demonstrating  that the zero-levelset of the differential zero-sum two-person game in an  HJ-Isaacs (HJI) setting~\citet{Isaacs1965,Evans1984} constitutes the safe set of a reachability problem~\citet{LygerosReachability}.  %We refer interested readers on the technicalities of the theory to Mitchell's paper~\citet{Mitchell2005}. In the rest of this paper, we focus on proximal approximators for computing reachable sets in the real of safety analysis of a dynamical system.
%
%The prevailing approach for computing the approximate solution to \eqref{eq:ivp-viscous} in reachability settings is via grid discretization and finite-difference schemes~\citet{CrandallTwoApprox} coupled with total variation diminishing (TVD) and  Courant–Friedrichs–Lewy (CFL) conditioning of the solution to assure numerical robustness~\citet{LevelSetsBook}. 
%It is well-known that solving these equations under appropriate boundary and/or initial conditions using the method of characteristics is limiting as a result of crossing characteristics~\citet{Crandall1983viscosity}. And global analysis is virtually impossible owing to the lack of existence and uniqueness of  solutions %$\lowervalue \in C^1(\Omega) \times (0, T]$ even if the Hamiltonian and boundary functions are smooth~\citet{Crandall1983viscosity}. This has motivated viscosity solutions~\citet{CrandallTwoApprox} that traverse the limit $\delta \rightarrow 0$. 
Popular numerically consistent optimal solutions~\citet{CrandallTwoApprox, CrandallMethodFracSteps} to \eqref{eq:ivp-viscous} utilize problem space grid discretization~\citet{LevelSetsBook, Mitchell2005}. However, these exact an exponential memory storage --- making their practical applications on large problems limited. Recent efforts~\citet{MoluxLevPyCDC, LevelSetTOMS} leveraged single-instruction, multiple threads (SIMT) with modern GPUs to accelerate computation on a host of practical problems --- showing improved spatial scalability with up to $76-92\%$ computational improvement on higher-dimensional problems. However, this still employed grid-storage of the value function; as problem dimensions increase, the need for more memory renders such approaches unattractive. This work thus tack another thread, jettisoning grid-discretization methods for sampling-based proximal operators in a statistical moving average context. This resolves the global extrema of the Moreau envelopes of HJ cost functions, which can be used to recover the reachable set  in nonconvex optimization settings~\citet{HJMAD, HJProx, Chaudhari2018, ChaudhariEntropy}. 

The Cole-Hopf transformation~\citet{EvansPDEBook, Chaudhari2018} employed in PDE analysis relates the heat equation's solution to that of \eqref{eq:ivp-viscous}. The essence of this work involves
%
\begin{inparaenum}[(i)]
	\item first, leveraging the Cole-Hopf transformation to linearize the nonlinear HJ equation\eqref{eq:ivp-viscous} to a heat equation. \label{contrib:cole-hopf}
	%
	\item The computation of the spatial gradients of $\valuefunc^\delta$ is a computationally intensive one in numerical HJ settings that typically employs upwinding schmes~\citet{LevelSetsBook}, essentially non-oscillatory total variation diminishing Runge-Kutta~\citet{OsherShuENO} and Courants-Friedrichs-Lax adaptive time step weighting to recover a smooth solution to $\partial_{\state}\valuefunc^\delta$. \label{contrib:upwinding}
	%
	\item We  therefore device a proximal operator equivalent to \ref{contrib:upwinding} that enables a sampling-based evaluation of $\partial_{\state}\valuefunc^\delta$ in a stochastic and adaptive Moreau gradient descent algorithmic scheme under two inputs in a differential game setting.  
	%
	\item The solution $\valuefunc^\delta$ to \eqref{eq:ivp-viscous} is then recovered as a logarithm of the solution to the surrogate heat equation (devised via the Cole-Hopf transformation \ref{contrib:cole-hopf}). 
	%
	\item Recovering the \textit{robustly controlled backward reachable set or tube}(RCBRT)~\citet{Mitchell2020} in the terminal Hamilton-Jacobi-Isaacs equation~\citet{Isaacs1965, Mitchell2005, MoluxLevPyCDC, LevelSetTOMS} of a typical reachability problem is thereafter turned into a lookup search problem for where the states fulfill the RCBRT condition. %Similar to our previous computational efforts, we accelerate numerical solutions on modern GPU hardware and provide open-source code for would-be practitioners. Ours is a   first-order proximal operator and Moreau envelope~\citet{Moreau} nonconvex computational optimization tool for resolving the ``generalized" viscosity solutions to HJ-Isaacs (HJI) PDE in  reachability computational settings.	
\end{inparaenum} 

The rest of this paper is structured as follows: \S \ref{sec:back}, sets the background for this work. Section \ref{sec:methods} describes the theoretical and algorithmic machinery that sets the numerical results of section \ref{sec:results} in order. The paper is concluded in \S \ref{sec:conclude}. 